\section{Method}
\label{sec:method}

\subsection{Environment}

We use a custom 3-role cooperative environment built on MiniGrid~\cite{minigrid}, consisting of a 15$\times$15 grid with three capture zones (top-left, top-right, center) and two moving enemies. Three agents are assigned one of three fixed roles: \textbf{Strike} (ranged DPS), \textbf{Vanguard} (melee tank), and \textbf{Support} (healer). Episodes run for 200 steps. The win condition requires the team to capture and hold 2 of 3 zones for 40 consecutive steps.

\subsection{Entity-Centric Observation}

Each agent receives a 63-dimensional entity-centric observation structured as four named token groups:

\begin{table}[h]
\centering
\caption{Observation layout (63 dimensions). Role one-hot uses [1,0,0]=Strike, [0,1,0]=Vanguard, [0,0,1]=Support.}
\label{tab:obs}
\begin{tabular}{llcp{6cm}}
\toprule
Token & Dims & Size & Content \\
\midrule
\texttt{self\_tok}   & [0:8]   & 8  & pos(2) + vel(2) + health(1) + role\_onehot(3) \\
\texttt{ally\_tok}   & [8:24]  & 16 & 2 allies $\times$ 8: rel\_pos(2) + health(1) + alive(1) + role\_onehot(3) + padding(1) \\
\texttt{enemy\_tok}  & [24:40] & 16 & 2 enemies $\times$ 8: rel\_pos(2) + health(1) + alive(1) + class\_conf(3) + padding(1) \\
\texttt{zone\_tok}   & [40:58] & 18 & 3 zones $\times$ 6: rel\_pos(2) + occupancy(1) + progress(1) + assigned(1) + value(1) \\
\texttt{ally\_mask}  & [58:60] & 2  & 0=valid, 1=padded (per ally slot) \\
\texttt{enemy\_mask} & [60:62] & 2  & 0=valid, 1=padded \\
\texttt{zone\_mask}  & [62:63] & 1  & 0=valid, 1=padded (zone 3 padding indicator) \\
\bottomrule
\end{tabular}
\end{table}

\textbf{Critical design choice:} The role one-hot appears in both \texttt{self\_tok} and each \texttt{ally\_tok} slot. This dual placement is load-bearing for the theoretical mechanism: it allows role labels to condition both the Query (from \texttt{self\_tok}) and the Key/Value (from \texttt{ally\_tok}) projections in cross-attention, enabling role-to-role relational routing.

\subsection{Reward Functions}

Role-specific rewards are mixed with a team component at $\alpha=0.2$:
\[
r_{\text{final}} = 0.8 \cdot r_{\text{individual}} + 0.2 \cdot \bar{r}_{\text{team}}
\]

Individual reward components per role:

\begin{table}[h]
\centering
\caption{Role-specific reward shaping.}
\label{tab:rewards}
\begin{tabular}{lp{9cm}}
\toprule
Role & Reward components \\
\midrule
\textbf{Strike}   & Approach bonus (scaled by distance to nearest zone) + capture bonus + too-close penalty ($-$0.1 when dist $<$ 2 to enemy) \\
\textbf{Vanguard} & Approach bonus + capture bonus + melee bonus ($+$0.1 per enemy within dist $<$ 2 in zone) \\
\textbf{Support}  & Approach-to-lowest-HP-ally bonus + heal bonus + boost bonus ($+$0.5 per nearby occupied zone within radius 3.0) \\
\bottomrule
\end{tabular}
\end{table}

\subsection{Policy Architectures}

All five conditions share the same observation format and action space (6 discrete actions). They differ in architectural structure and whether role labels are present.

\subsubsection{Condition 1 — Full (Proposed)}

The full architecture processes each entity set through a two-stage attention pipeline:

\paragraph{Stage 1: Intra-Set Self-Attention.}
Each entity set is first encoded by a linear projection ($\mathbb{R}^{d_{\text{in}}} \to \mathbb{R}^{64}$), then processed by an intra-set multi-head self-attention block:
\begin{align*}
e_{\text{ally}}^{(i)} &= \text{Linear}(8, 64)(x_{\text{ally}}^{(i)}) \\
\tilde{e}_{\text{ally}} &= \text{LayerNorm}\left(e_{\text{ally}} + \text{MHA}_{\text{ally}}(e_{\text{ally}}, e_{\text{ally}}, e_{\text{ally}})\right)
\end{align*}
Analogous blocks are applied to enemy and zone sets. Padding masks prevent attention over invalid slots; a \texttt{\_safe\_mask()} guard forces-unmasks slot 0 when all slots are padded to prevent NaN.

\paragraph{Stage 2: Cross-Attention.}
The encoded \texttt{self\_tok} embedding serves as the Query for three separate cross-attention heads\textemdash one per entity set (ally, enemy, zone):
\[
\text{ctx}_{\text{ally}} = \text{MHA}_{\text{cross-ally}}(e_{\text{self}},\, \tilde{e}_{\text{ally}},\, \tilde{e}_{\text{ally}})
\]

\paragraph{Output MLP.}
The three context vectors and the self embedding are concatenated and passed through a two-layer MLP:
\[
a = \text{Linear}(128, |\mathcal{A}|)\left(\text{ReLU}\left(\text{Linear}(64 \cdot 3 + 64,\, 128)\left([\text{ctx}_{\text{ally}};\, \text{ctx}_{\text{enemy}};\, \text{ctx}_{\text{zone}};\, e_{\text{self}}]\right)\right)\right)
\]

\subsubsection{Condition 2 — No Self-Attention}

Same as Condition~1, but the three intra-set self-attention blocks are removed. Encoded embeddings feed directly into cross-attention without intra-set contextualization. This ablation isolates the contribution of intra-set relational structure.

\subsubsection{Condition 3 — Role Label + Pure MLP}

All attention is removed. The full 63-dimensional observation is flattened and processed by a three-layer MLP:
\[
a = \text{Linear}(128, |\mathcal{A}|)\left(\text{ReLU}\left(\text{Linear}(256, 128)\left(\text{ReLU}\left(\text{Linear}(63, 256)(o)\right)\right)\right)\right)
\]
Role labels remain present in the input. This is the primary comparison condition for the structural routing claim.

\subsubsection{Condition 4 — No Role Label (Shared Policy)}

Same architecture as Condition~1, but role one-hot dimensions are zeroed in both \texttt{self\_tok} and all \texttt{ally\_tok} slots before any encoding. A single shared policy network handles all three roles. Role-specific rewards still apply during training. This condition isolates the contribution of role label information: with independent policies, role differentiation arises trivially from trajectory separation; only a shared policy makes label removal a meaningful test.

\subsubsection{Condition 5 — Flat MLP Baseline}

Same as Condition~3 but role one-hot dimensions are also zeroed. Pure flat MLP, no labels, no attention.

\subsection{Training Setup}

We use Multi-Agent PPO (MAPPO)~\cite{yu2022mappo} with a centralized critic taking the concatenated observations of all agents ($63 \times 3 = 189$ dims). All hyperparameters are fixed across all five conditions:

\begin{table}[h]
\centering
\caption{Training hyperparameters (identical across all conditions).}
\label{tab:hparams}
\begin{tabular}{lc}
\toprule
Hyperparameter & Value \\
\midrule
Algorithm         & MAPPO \\
Total steps       & 2,000,000 \\
Batch size        & 512 \\
Mini-batch size   & 64 \\
PPO epochs/update & 4 \\
Discount $\gamma$ & 0.99 \\
GAE $\lambda$     & 0.95 \\
PPO clip $\varepsilon$ & 0.2 \\
Entropy coefficient & 0.01 \\
Learning rate     & $3 \times 10^{-4}$ (Adam) \\
Max grad norm     & 0.5 \\
Seeds             & 3 \\
\bottomrule
\end{tabular}
\end{table}

Conditions 1, 2, 3, and 5 use three independent role-specific policies; each network is updated exclusively on trajectories from its assigned role. Condition~4 uses a single shared policy updated on all trajectories. Checkpoints are saved every 200k steps; win rate is evaluated over 10 episodes every 10k steps.
